\documentclass[
  aspectratio=169,
  font=times,
  sectionpage=true,
]{tumbeamer}

\usepackage{booktabs}
\usepackage{array}
\usepackage{tabularx}

\usepackage{tikz}
\usetikzlibrary{arrows.meta, positioning}
\usepackage{pgfplots}
\pgfplotsset{compat=1.18}

\newcolumntype{L}[1]{>{\raggedright\arraybackslash}p{#1}}

\usepackage{etoolbox}

% Reusable five-stage agent-loop diagram.
% Optional arg = stage to highlight (intent|plan|gen|exec|refine); default none.
\newcommand{\agentloopdiagram}[1][none]{%
  \def\hlintent{stage}\def\hlplan{stage}\def\hlgen{stage}%
  \def\hlexec{stage}\def\hlrefine{stage}%
  \ifstrequal{#1}{intent}{\def\hlintent{here}}{}%
  \ifstrequal{#1}{plan}{\def\hlplan{here}}{}%
  \ifstrequal{#1}{gen}{\def\hlgen{here}}{}%
  \ifstrequal{#1}{exec}{\def\hlexec{here}}{}%
  \ifstrequal{#1}{refine}{\def\hlrefine{here}}{}%
  \begin{tikzpicture}[
    stage/.style={draw=TUMBlue, fill=TUMBlue!8, rounded corners=3pt,
                  font=\small, align=center, minimum height=0.8cm,
                  minimum width=2.2cm},
    here/.style={draw=TUMOrange, fill=TUMOrange!15, rounded corners=3pt,
                 font=\small\bfseries, align=center, minimum height=0.8cm,
                 minimum width=2.2cm},
    note/.style={font=\scriptsize, text=TUMGrayDark, align=center, inner sep=1pt},
    arr/.style={-{Stealth[length=1.8mm]}, very thick, TUMGrayDark},
    feedback/.style={-{Stealth[length=1.8mm]}, very thick, dashed, TUMOrange},
    newtask/.style={-{Stealth[length=1.4mm]}, thin, dotted, TUMGray},
    fblab/.style={font=\tiny, text=TUMOrange, align=center, inner sep=1.5pt,
                  fill=white},
    ntlab/.style={font=\tiny, text=TUMGrayDark, align=center, inner sep=1.5pt,
                  fill=white},
    key/.style={font=\scriptsize, text=TUMGrayDark, anchor=west},
  ]
    \node[\hlintent] (intent) at (0,0)      {Intent\\capture};
    \node[\hlplan]   (plan)   at (3.3,0)    {Planning};
    \node[\hlgen]    (gen)    at (6.6,0)    {Code\\generation};
    \node[\hlexec]   (exec)   at (6.6,-1.9) {Execution};
    \node[\hlrefine] (refine) at (3.3,-1.9) {Refinement};

    \draw[arr] (intent) -- (plan);
    \draw[arr] (plan) -- (gen);
    \draw[arr] (gen) -- (exec);
    \draw[arr] (exec) -- (refine);

    \draw[newtask] (refine.west) to[out=180, in=270]
      node[ntlab, pos=0.55] {new task,\\not feedback} (intent.south);

    \draw[feedback] (refine.north) -- node[fblab, right] {replan} (plan.south);
    \draw[feedback] (refine.north east) -- node[fblab, pos=0.5] {regenerate}
      (gen.south west);

    \node[note, above=0.15cm of intent] {NL issue or spec};
    \node[note, below=0.15cm of exec] {test failures, traces};

    % Key: which visual element each cross-cutting control acts on
    \draw[arr] (0,-3.15) -- ++(0.55,0);
    \node[key] at (0.78,-3.15)
      {\textbf{Orchestration} --- the arrows: which step runs next};
    \node[stage, minimum width=0.55cm, minimum height=0.34cm] at (0.27,-3.72) {};
    \node[key] at (0.78,-3.72)
      {\textbf{Interface \& context} --- inside each stage};
  \end{tikzpicture}%
}

\TUMbeamersetup{
  title page   = TUM tower,
  content page = TUM more space,
  footline     = TUM infoline,
  footline on  = {content page},
}

\makeatletter
\AtBeginDocument{%
  \setbeamertemplate{footline}[TUM infoline]%
  \beamer@calculateheadfoot
}
\makeatother

\makeatletter
\AtBeginSection[]{%
  \TUMframestyle{section page}%
  {\setbeamertemplate{footline}{}%
   \begin{frame}[c, noframenumbering]
     \sectionpage
   \end{frame}}%
  \TUMframestyle{content page}%
}
\makeatother

\setlength{\leftmargini}{1.4em}

\title{Autonomous Coding Agents}
\subtitle{A seminar survey of agentic AI \\ for code generation --- SS26}
\author{Thua Duc Nguyen}
\institute{Supervisor: Derui Zhu \\ \theDepartmentName, \theUniversityName}
\date{\today}

\begin{document}

\maketitle

\begin{frame}{How we write code has changed}
  \small
  \begin{tabular}{@{}L{0.14\textwidth}L{0.26\textwidth}L{0.24\textwidth}L{0.24\textwidth}@{}}
    \toprule
      & \textbf{Machine} & \textbf{Human} & \textbf{Checked by} \\
    \midrule
    Manual
      & Compiles and runs it
      & Writes every line
      & The test suite \\
    \addlinespace
    Autocomplete \newline {\scriptsize 2021--}
      & Suggests the next lines
      & Accepts or rejects each
      & You, line by line \\
    \addlinespace
    One-shot \newline {\scriptsize 2022--}
      & Writes a whole function
      & Specifies, reads all of it
      & \textbf{Nothing} \\
    \addlinespace
    \textbf{Agentic} \newline {\scriptsize 2024--}
      & Finds files, edits, runs tests
      & States intent, reviews outcome
      & Tests, \emph{during} generation \\
    \bottomrule
  \end{tabular}

  \medskip
  \begin{center}
    \begin{tikzpicture}[
      unit/.style={draw=TUMBlue, fill=TUMBlue!8, rounded corners=2pt,
                   font=\small, align=center, minimum height=6.5mm},
      arr/.style={-{Stealth[length=2mm]}, thick, TUMGrayDark},
      era/.style={font=\scriptsize, text=TUMGrayDark, align=center, inner sep=1pt},
      reads/.style={font=\scriptsize, text=TUMGrayDark, align=center, inner sep=1pt},
    ]
      \node[unit, minimum width=15mm] (u1) at (0,0)   {a line};
      \node[unit, minimum width=27mm] (u2) at (3.1,0) {a function};
      \node[unit, minimum width=41mm] (u3) at (7.1,0) {a whole issue};

      \draw[arr] (u1) -- (u2);
      \draw[arr] (u2) -- (u3);

      \node[era, below=0.8mm of u1] (e1) {2021--};
      \node[era, below=0.8mm of u2] (e2) {2022--};
      \node[era, below=0.8mm of u3] (e3) {2024--};

    \end{tikzpicture}
  \end{center}
\end{frame}

\begin{frame}{Real progress}
  \medskip
  \textbf{SWE-bench Verified} progress --- 500 human-screened GitHub issues, graded by the project's own tests.

  \smallskip
  \begin{center}
    \begin{tikzpicture}
      \begin{axis}[
        width=0.88\textwidth,
        height=3.8cm,
        ylabel={\% resolved},
        ymin=0, ymax=104,
        ytick={0,20,40,60,80,100},
        xmin=2024.10, xmax=2026.66,
        xtick={2024.25, 2024.75, 2025.25, 2025.75, 2026.25},
        xticklabels={Apr'24, Oct'24, Apr'25, Oct'25, Apr'26},
        xticklabel style={font=\scriptsize},
        yticklabel style={font=\scriptsize},
        ylabel style={font=\scriptsize},
        grid=major,
        grid style={dashed, TUMGrayLight},
        axis line style={TUMGrayDark},
        tick style={TUMGrayDark},
        every axis plot/.append style={thick, mark size=1.6pt},
      ]
      \addplot[TUMBlue, mark=*] coordinates {
        (2024.251,18.2) (2024.467,33.6) (2024.571,23.2) (2024.918,50.8)
        (2024.918,55.0) (2024.975,57.2) (2025.044,64.6) (2025.099,63.4)
        (2025.203,65.4)
      };

      \addplot[TUMBlue, mark=*] coordinates {
        (2025.203,65.4) (2026.25,80.9)
      };

      \addplot[TUMOrange, mark=square*] coordinates {
        (2025.71,23.0) (2026.25,45.9)
      };

      \draw[<->, TUMGrayDark, thin] (axis cs:2026.25,79.0)
        -- node[font=\tiny, text=TUMGrayDark, right=1pt, inner sep=1pt]
           {$-35$\,pp} (axis cs:2026.25,47.8);

      \node[font=\scriptsize, text=TUMBlue, anchor=north west, inner sep=2pt]
        at (axis cs:2024.16,16.0) {SWE-agent 18.2\,\%};
      \node[font=\scriptsize, text=TUMBlue, anchor=south east, inner sep=2pt]
        at (axis cs:2025.15,68.0) {Augment 65.4\,\%};
      \node[font=\scriptsize, text=TUMBlue, anchor=south east, inner sep=2pt]
        at (axis cs:2026.36,83.0) {Opus 4.5, Verified 80.9\,\%};
      \node[font=\scriptsize, text=TUMOrange, anchor=north east, inner sep=2pt]
        at (axis cs:2026.36,43.5) {Pro 45.9\,\%};
      \node[font=\scriptsize, text=TUMOrange, anchor=north east, inner sep=2pt]
        at (axis cs:2025.60,19.0) {Pro at release $\sim$23\,\%};
      \end{axis}
    \end{tikzpicture}

    {\footnotesize Blue: \textbf{SWE-bench Verified.} \quad Orange:
    \textbf{SWE-bench Pro} --- harder, released Sept~2025.}
  \end{center}

  \smallskip
  \textbf{Massive improvement in two years}, and the models were
  not trained to be agents.

  \smallskip
  {\footnotesize More on SWE-bench later.}
\end{frame}

\begin{frame}{What this talk shows}
  \begin{enumerate}
    \item \textbf{How do the agentic AI systems work?} \\
      {\small Five stages, orchestration and interface, planning and
      refinement, five systems.}
      \medskip
    \item \textbf{How do we measure performance?} \\
      {\small The benchmark, the metrics, what the numbers hide.}
      \medskip
    \item \textbf{What are the challenges?} \\
      {\small Three open challenges --- and the pattern behind them.}
  \end{enumerate}
\end{frame}

\begin{frame}{A few terms, up front}
  \begin{tabular}{@{}L{0.20\textwidth}L{0.72\textwidth}@{}}
    \toprule
    \textbf{Scaffold} & The non-model code around it --- tools, prompts,
      control flow \\
    \addlinespace
    \textbf{AST} & A parsed structure of the code --- classes, functions,
      calls, not raw text \\
    \addlinespace
    \textbf{Harness} & The environment a result is run and graded in \\
    \addlinespace
    \textbf{Oracle} & Something outside the model that can check it --- a
      test suite, a compiler, a human \\
    \addlinespace
    \textbf{Critic} & A separate model, trained to score candidate outputs \\
    \bottomrule
  \end{tabular}
\end{frame}

\section{How do the agentic AI systems work?}

\begin{frame}{What is an agentic coding system?}
  A system that, given a natural-language request, \textbf{autonomously} produces a
  working code change --- navigating files, editing, executing, and revising
  without step-by-step human guidance.

  \bigskip
  \begin{itemize}
    \item \textbf{Autonomous} --- decides what to do next, not told
      \medskip
    \item \textbf{Grounded in execution} --- runs tests, reads errors, acts on
      them
      \medskip
    \item \textbf{Repository-scale} --- works across files in a real codebase,
      not on isolated functions
  \end{itemize}
\end{frame}

\begin{frame}{The agent loop}
  \label{fr:stages}
  \begin{center}
    \agentloopdiagram
  \end{center}
\end{frame}

\begin{frame}{How the loop is controlled}
  The five stages are \textbf{what} happens. Two controls decide \textbf{how} the
  loop runs --- and each acts on a different part of it:

  \medskip
  \begin{center}
    \begin{tabular}{@{}lll@{}}
      \toprule
       & \textbf{Orchestration} & \textbf{Interface \& context} \\
      \midrule
      Acts on & the \textbf{arrows} between stages & \textbf{inside} each stage \\
      \addlinespace
      What it does & decides the next step & sets what it can see and do \\
      \bottomrule
    \end{tabular}
  \end{center}

  \medskip
  Neither is a stage --- both cut \emph{across} all five, and decide whether the
  loop even closes.
\end{frame}

\begin{frame}{Orchestration}
  \label{fr:orch}
  Two common ways to run the same five stages.

  \bigskip

  \begin{tabular}{@{}L{0.22\textwidth}L{0.46\textwidth}L{0.22\textwidth}@{}}
    \toprule
    \textbf{Mode} & \textbf{How it works} & \textbf{Who decides} \\
    \midrule
    Agentic loop \newline {\scriptsize ReAct}~\cite{yao2023react}
      & Thought, Action, Observation; repeat
      & The model \\
    \addlinespace
    Fixed pipeline~\cite{xia2024agentless}
      & Localise, repair, validate; no loop
      & The scaffold \\
    \bottomrule
  \end{tabular}
\end{frame}

\begin{frame}{Orchestration: agentic loop}
  \label{fr:orch-agentic}
  {\small The \textbf{model} is in control: the \textbf{ReAct} loop --- Thought,
  Action, Observation --- repeated until \emph{the model} decides it is
  done~\cite{yao2023react}.}

  \bigskip
  \begin{center}
    \begin{tikzpicture}[
      stage/.style={draw=TUMBlue, fill=TUMBlue!8, rounded corners=3pt,
                    font=\small, align=center, minimum height=0.8cm,
                    minimum width=2.4cm},
      arr/.style={-{Stealth[length=1.8mm]}, very thick, TUMGrayDark},
      loop/.style={-{Stealth[length=1.8mm]}, very thick, dashed, TUMOrange},
    ]
      \node[stage] (think) at (0,0)   {Thought};
      \node[stage] (act)   at (3.2,0) {Action\\{\scriptsize tool call}};
      \node[stage] (obs)   at (6.4,0) {Observation};
      \draw[arr] (think) -- (act);
      \draw[arr] (act) -- (obs);
      \draw[loop] (obs.south) to[out=270, in=270]
        node[below=1mm, font=\scriptsize, text=TUMOrange]
        {model loops until done} (think.south);
    \end{tikzpicture}
  \end{center}

  \bigskip
  {\small Read a file, run tests, edit --- \textbf{any tool, any order,} as many
  rounds as it takes. Nothing outside the model fixes the sequence.}
\end{frame}

\begin{frame}{Orchestration: fixed pipeline}
  \label{fr:orch-fixed}
  {\small The \textbf{scaffold} is in control: the steps are \textbf{fixed in
  advance}; the model only fills each one~\cite{xia2024agentless}.}

  \bigskip
  \begin{center}
    \begin{tikzpicture}[
      stage/.style={draw=TUMBlue, fill=TUMBlue!8, rounded corners=3pt,
                    font=\small, align=center, minimum height=0.8cm,
                    minimum width=2.2cm},
      stop/.style={draw=TUMGray, fill=none, rounded corners=3pt,
                   font=\small, align=center, minimum height=0.8cm,
                   minimum width=1.6cm, text=TUMGrayDark},
      arr/.style={-{Stealth[length=1.8mm]}, very thick, TUMGrayDark},
    ]
      \node[stage] (loc) at (0,0)   {Localise};
      \node[stage] (rep) at (2.9,0) {Repair};
      \node[stage] (val) at (5.8,0) {Validate};
      \node[stop]  (end) at (8.5,0) {stop};
      \draw[arr] (loc) -- (rep);
      \draw[arr] (rep) -- (val);
      \draw[arr] (val) -- (end);
    \end{tikzpicture}
  \end{center}

  \bigskip
  {\small On failure it stops or moves on --- \textbf{no arrow goes back.}
  Cheaper and predictable, but no recovery.}
\end{frame}

\begin{frame}{Interface and Context}

  What the agent sees and does.

  \bigskip

  \textbf{Agent--Computer Interface (ACI)}
  \begin{itemize}
    \item Purpose-built command set
    \item Scrollable file viewer, linter-guarded edits, summarised search
    \item The scaffold controls what the model can see
    \item Malformed actions rejected before they run
  \end{itemize}

  \medskip
  \textbf{Code as action}
  \begin{itemize}
    \item The agent writes Python or bash in a sandbox
    \item Full shell access, no guardrails, no summarisation
    \item The model manages its own context
  \end{itemize}
\end{frame}

\begin{frame}{Interface and Context}
  \label{fr:iface}
  \small
  \begin{tabular}{@{}L{0.14\textwidth}L{0.40\textwidth}L{0.36\textwidth}@{}}
    \toprule
    & \textbf{ACI}~\cite{yang2024sweagent}
    & \textbf{Code as action}~\cite{wang2024codeact} \\
    \midrule
    \textbf{View files}
      & \texttt{open file.py 100} \newline {\scriptsize 100-line window,
        scrollable}
      & \texttt{cat file.py} \newline {\scriptsize raw output, entire file} \\
    \addlinespace
    \textbf{Edit}
      & \texttt{edit 42:45 <new code>} \newline {\scriptsize linter rejects
        syntax errors}
      & \texttt{sed -i ...} or write Python \newline {\scriptsize no guard} \\
    \addlinespace
    \textbf{Search}
      & \texttt{search\_dir "pattern"} \newline {\scriptsize summarised
        matches}
      & \texttt{grep -rn "pattern" .} \newline {\scriptsize raw hits} \\
    \addlinespace
    \textbf{Execute}
      & \texttt{python -m pytest} \newline {\scriptsize via tool call}
      & \texttt{bash: pytest} \newline {\scriptsize direct shell} \\
    \bottomrule
  \end{tabular}
\end{frame}

\begin{frame}{Back to the five stages}
  \begin{center}
    \scalebox{0.9}{\agentloopdiagram[plan]}
  \end{center}

  \medskip
  Now, stage by stage --- starting with \textbf{Planning}, where the model has
  the \textbf{least to lean on}.
\end{frame}

\begin{frame}{Plan: finding the right code}
  \label{fr:planning}
  Two commitments: \textbf{where} to change the code --- which file, which
  function --- and \textbf{what} the change should be. Both are guesses; both
  need evidence that could confirm them.

  \bigskip
  \begin{itemize}
    \item \textbf{Where} --- among tens of thousands of functions, which one?
      The relevant symbol often does not appear in the issue text.
      \medskip
    \item \textbf{What} --- a signature, an interface, an invariant to hold ---
      stated concretely enough that the repository can check it.
  \end{itemize}

\end{frame}

\begin{frame}{Plan: shape}
  \label{fr:shape}
  {\large How does the agent \textbf{explore possible plans}?}

  \begin{center}
  \scalebox{0.85}{%
  \begin{tikzpicture}[
    node distance=0.8cm and 1.2cm,
    thought/.style={draw, circle, fill=TUMBlue!10, minimum size=0.6cm,
                    font=\tiny},
    arr/.style={-{Stealth[length=1.2mm]}, semithick},
    backarr/.style={-{Stealth[length=1.2mm]}, semithick, dashed},
  ]
    \node[thought] (c1) {A};
    \node[thought, right=of c1] (c2) {B};
    \node[thought, right=of c2] (c3) {C};
    \node[thought, right=of c3] (c4) {D};
    \draw[arr] (c1) -- (c2);
    \draw[arr] (c2) -- (c3);
    \draw[arr] (c3) -- (c4);
    \node[above=0.15cm of c2, font=\scriptsize\bfseries]
      {Chain of Thought~\cite{wei2022chain}};

    \node[thought, below=1.15cm of c1] (t1) {A};
    \node[thought, below right=0.5cm and 0.4cm of t1] (t2a) {B$_2$};
    \node[thought, below left=0.5cm and 0.4cm of t1] (t2b) {B$_1$};
    \node[thought, below right=0.5cm and 0.2cm of t2a] (t3a) {C$_3$};
    \node[thought, below left=0.5cm and 0.2cm of t2a] (t3b) {C$_2$};
    \node[thought, below=0.7cm of t2b] (t3c) {C$_1$};
    \draw[arr] (t1) -- (t2a);
    \draw[arr] (t1) -- (t2b);
    \draw[arr] (t2a) -- (t3a);
    \draw[arr] (t2a) -- (t3b);
    \draw[arr] (t2b) -- (t3c);
    \draw[backarr, TUMOrange] (t3b) to[bend right=45] (t2a);
    \node[above=0.15cm of t1, font=\scriptsize\bfseries]
      {Tree of Thoughts~\cite{yao2023tree}};
    \node[below=0.05cm of t3b, font=\tiny, text=TUMOrange] {backtrack};

    \node[thought, below=1.15cm of c3] (g1) {A};
    \node[thought, below right=0.5cm and 0.4cm of g1] (g2a) {B$_2$};
    \node[thought, below left=0.5cm and 0.4cm of g1] (g2b) {B$_1$};
    \node[thought, below=1.2cm of g1] (g3) {C};
    \node[thought, below=0.7cm of g3] (g4) {D};
    \draw[arr] (g1) -- (g2a);
    \draw[arr] (g1) -- (g2b);
    \draw[arr, TUMBlue] (g2a) -- (g3);
    \draw[arr, TUMBlue] (g2b) -- (g3);
    \draw[arr] (g3) to[bend right=40]
      node[midway, right, font=\tiny, inner sep=1pt] {revisit} (g2a);
    \draw[arr] (g3) -- (g4);
    \node[above=0.15cm of g1, font=\scriptsize\bfseries]
      {Graph of Thoughts~\cite{besta2024graph}};
    \node[left=0.05cm of g3, font=\tiny, text=TUMBlue] {merge};
  \end{tikzpicture}}
  \end{center}

  \vspace{-2mm}
  \centering
  {\small \textbf{But which path is right?} Nothing has run yet, so the model is
  only \textbf{guessing}.}
\end{frame}

\begin{frame}{Plan: Chain of Thought}
  \label{fr:bu-cot}
  \textbf{The problem to plan for} --- no file, no fix given:

  \smallskip
  \begin{center}
    \begin{tikzpicture}
      \node[draw, rounded corners, fill=TUMOrange!8, inner sep=6pt,
            text width=0.86\textwidth, font=\small, align=left]
        {\textit{``The app crashes when you check out with an empty cart.''}};
    \end{tikzpicture}
  \end{center}

  \medskip
  \begin{center}
  \begin{tikzpicture}
    \node[draw, rounded corners, fill=TUMBlue!5, inner sep=8pt,
          text width=0.9\textwidth, font=\ttfamily\small, align=left]
      {\textcolor{TUMBlue}{Thought:} empty cart --- maybe a divide-by-zero in the total.\\
       \textcolor{TUMOrange}{Action:} grep 'def total'\\
       \textcolor{black!55}{Obs:} cart.py:12 --- return sum(prices) / len(items)\\[3pt]
       \textcolor{TUMBlue}{Thought:} len(items) is 0 when the cart is empty.\\
       \textcolor{TUMOrange}{Action:} open cart.py:12 --- confirm the division\\
       \textcolor{black!55}{Obs:} yes, no empty-case guard\\[4pt]
       \textbf{Plan} --- \emph{where:} cart.py:12 \texttt{total()}\quad
       \emph{what:} guard empty cart $\to$ 0};
  \end{tikzpicture}
  \end{center}
\end{frame}

\begin{frame}{Plan: Tree of Thoughts}
  \label{fr:bu-tot}
  {\small No agent grows a tree of half-thoughts. The tree is \textbf{whole ReAct
  runs}, sampled and \textbf{scored by something outside the model}.}

  \medskip
  \begin{center}
  \begin{tikzpicture}[
    b/.style={draw, rounded corners, fill=TUMBlue!10, font=\scriptsize,
              inner sep=4pt},
    arr/.style={-{Stealth[length=1.4mm]}, semithick},
  ]
    \node[b] (i) at (0,0) {issue};
    \node[b] (r1) at (3.5,1.3)  {ReAct run 1 $\rightarrow$ patch A};
    \node[b] (r2) at (3.5,0)    {ReAct run 2 $\rightarrow$ patch B};
    \node[b] (r3) at (3.5,-1.3) {ReAct run 3 $\rightarrow$ patch C};
    \node[b] (s1) at (7.7,1.3)  {tests 2/3};
    \node[b, fill=TUMBlue!30] (s2) at (7.7,0) {tests 3/3 --- \textbf{pick}};
    \node[b] (s3) at (7.7,-1.3) {tests 0/3};
    \draw[arr] (i) -- (r1);
    \draw[arr] (i) -- (r2);
    \draw[arr] (i) -- (r3);
    \draw[arr] (r1) -- (s1);
    \draw[arr] (r2) -- (s2);
    \draw[arr] (r3) -- (s3);
  \end{tikzpicture}
  \end{center}

  \medskip
  {\small The scorer is \textbf{tests or a trained critic}, not the model's
  hunch. OpenHands \textbf{60\%\,$\rightarrow$\,66\%} with five
  reranked~\cite{openhands2025verified}; SWE-Search wraps it in MCTS with a
  learned value --- \textbf{+23\%}, but a model call per
  node~\cite{antoniades2024swesearch}.}
\end{frame}

\begin{frame}{Plan: Graph of Thoughts}
  \label{fr:bu-got}
  {\small Merging branches, revisiting nodes --- shown on \textbf{puzzles}
  (Game-of-24, sorting), \textbf{almost never on repository
  bugs}~\cite{besta2024graph}.}

  \medskip
  \begin{itemize}
    \item In real SWE agents you see the \textbf{two ends}: one chain (CoT), or
      many chains scored (best-of-$N$).
      \medskip
    \item The closest thing to \emph{revisit}: \textbf{Reflexion} writes a note
      after a failed attempt and retries with it --- a loose loop, not a
      graph~\cite{shinn2023reflexion}.
      \medskip
    \item No surveyed end-to-end SWE agent builds a literal thought-graph.
  \end{itemize}

  \medskip
  \begin{block}{Takeaway}
    The graph is the \emph{conceptual} top of the spectrum --- useful to reason
    about, not something you will find running in production.
  \end{block}
\end{frame}

\begin{frame}{Generate: turning it into code}
  The model writes or edits source code --- a patch, a new function, a refactored
  block.

  \medskip
  \begin{itemize}
    \item \textbf{Infilling}
      \begin{itemize}
        \item completes a span from the surrounding context
        \item works well within a single function
      \end{itemize}
      \medskip
    \item \textbf{Diff generation}
      \begin{itemize}
        \item a unified diff, not a whole file --- smaller output
        \item fragile if the surrounding context is stale
      \end{itemize}
      \medskip
    \item \textbf{Sample $N$, pick one}
      \begin{itemize}
        \item many candidates, ranked by test results or a critic
        \item AlphaCode samples up to \emph{millions}, clusters by
          execution~\cite{li2022alphacode}
      \end{itemize}
  \end{itemize}
\end{frame}

\begin{frame}{Execute: running the code}
  The agent runs what it wrote. This is where the \textbf{environment can say no}.

  \medskip
  \begin{itemize}
    \item \textbf{Test suite}
      \begin{itemize}
        \item fail-to-pass tests must flip
        \item pass-to-pass must not regress
      \end{itemize}
      \medskip
    \item \textbf{Linter / type checker}
      \begin{itemize}
        \item catches syntax errors and type mismatches before tests run
      \end{itemize}
      \medskip
    \item \textbf{Error traces}
      \begin{itemize}
        \item stack traces and assertion messages say \emph{what} failed and
          \emph{where}
      \end{itemize}
  \end{itemize}

  \medskip
  {\small Right or wrong, that verdict is what refinement works with next.}
\end{frame}

\begin{frame}{Refinement: iterative self-repair}
  \label{fr:refine1}
  \small
  Execution returned a verdict. What the agent does with it is refinement.

  \medskip
  \begin{tabular}{@{}L{0.26\textwidth}L{0.30\textwidth}L{0.34\textwidth}@{}}
    \toprule
    \textbf{Mechanism} & \textbf{What closes the loop}
      & \textbf{Characteristic failure} \\
    \midrule
    Self-Refine~\cite{madaan2023selfrefine}
      & The model critiques itself
      & Plateaus \\
    \addlinespace
    Reflexion~\cite{shinn2023reflexion}
      & Verbal critique across attempts; self-written tests
      & Inherits what the tests missed \\
    \addlinespace
    Self-debugging~\cite{chen2024selfdebug}
      & Explains its own code; test verdict if any
      & Explains the bug as intended \\
    \bottomrule
  \end{tabular}

  \bigskip
  {\small Closed by the generator re-reading itself. Three that bring in
  something else --- next.}
\end{frame}

\begin{frame}{Refinement: tests and trained critics}
  \label{fr:refine2}
  \small
  \begin{tabular}{@{}L{0.26\textwidth}L{0.30\textwidth}L{0.34\textwidth}@{}}
    \toprule
    \textbf{Mechanism} & \textbf{What closes the loop}
      & \textbf{Characteristic failure} \\
    \midrule
    D4C~\cite{xu2024d4c}
      & Failing tests $+$ the whole function
      & Plausible but incorrect patches \\
    \addlinespace
    Agentless validate~\cite{xia2024agentless}
      & Regression $+$ generated tests
      & The generated half inherits the misreading \\
    \addlinespace
    Best-of-$N$ $+$ critic~\cite{openhands2025verified}
      & Executed rollouts, ranked by a trained critic
      & Cost scales with $N$ \\
    \bottomrule
  \end{tabular}

  \bigskip
  {\small Per-instance sandbox --- for \textbf{reproducibility} first,
  containment second.}
\end{frame}

\begin{frame}{Refinement: when does the loop stop?}
  \label{fr:refinelast}
  \begin{itemize}
    \item \textbf{All tests pass}
      \begin{itemize}
        \item but that is quietly an oracle
        \item the stopping rule has smuggled in ground truth
      \end{itemize}
      \medskip
    \item \textbf{Diminishing returns}
      \begin{itemize}
        \item at fixed budget, another repair round is often \emph{worse} than
          a fresh sample~\cite{olausson2024selfrepair}
      \end{itemize}
      \medskip
    \item \textbf{Iterate or best-of-$N$ --- the same lever}
      \begin{itemize}
        \item fix the budget to compare
      \end{itemize}
  \end{itemize}
\end{frame}

\begin{frame}{Five systems compared}
  \label{fr:systems}
  \scriptsize
  \begin{tabular}{@{}L{0.16\textwidth}L{0.14\textwidth}L{0.16\textwidth}L{0.18\textwidth}L{0.18\textwidth}@{}}
    \toprule
    \textbf{System}
      & \textbf{Orchestration}
      & \textbf{Interface}
      & \textbf{Planning}
      & \textbf{Refinement} \\
    \midrule
    SWE-agent
      & Agentic loop
      & ACI
      & Implicit (model)
      & Self-repair \\
    \addlinespace
    AutoCodeRover
      & Agentic loop
      & AST APIs
      & AST localisation
      & Re-generate \\
    \addlinespace
    OpenHands
      & Agentic loop
      & Code as action
      & Implicit (model)
      & Best-of-$N$ + critic \\
    \addlinespace
    Agentless
      & Fixed pipeline
      & Skeleton + emb.
      & Hierarchical loc.
      & Regression + gen.\ tests \\
    \addlinespace
    Devin
      & Agentic loop
      & Own VM
      & Knowledge store
      & --- \\
    \bottomrule
  \end{tabular}
\end{frame}

\section{How do we measure performance?}

\begin{frame}{SWE-bench}
  \label{fr:swebench}
  Merged pull requests that \textbf{close an issue} and \textbf{touch the test
  suite}~\cite{jimenez2024swebench}.

  \bigskip
  \begin{enumerate}
    \item \textbf{Given} --- the issue text $+$ the repository \emph{before}
      the merge.
    \item \textbf{Produce} --- a patch.
    \item \textbf{Graded} --- \emph{fail-to-pass} must flip,
      \emph{pass-to-pass} must not regress.
  \end{enumerate}
\end{frame}

\begin{frame}{Measuring generated code}
  \medskip
  \begin{itemize}
    \item \textbf{Functional correctness} --- does it run and pass? The only
      thing that counts.
      \smallskip
    \item \textbf{pass@$k$} --- give it $k$ tries; solved if any one
      passes~\cite{chen2021humaneval}.
      \smallskip
    \item \textbf{Resolve rate} --- pass@1 on a whole repo: solved on the first
      try.
      \smallskip
    \item \textbf{Cost} --- dollars, calls, time. A higher score can cost much
      more~\cite{sweeffi2025}.
  \end{itemize}
\end{frame}

\begin{frame}{SWE-bench: the five splits}
  \label{fr:splits}
  \begin{tabularx}{\textwidth}{@{}lll>{\raggedright\arraybackslash}X@{}}
    \toprule
    \textbf{Split} & \textbf{Size} & \textbf{Released} & \textbf{Selected how} \\
    \midrule
    Full & 2{,}294 & Oct 2023 & Every scraped PR that touches tests \\
    Lite & 300 & Mar 2024 & Single-file patches; some issues leak the
      fix~\cite{xia2024agentless} \\
    Verified & 500 & Aug 2024 & Human-screened for clear issues and fair
      tests~\cite{openai2024verified} \\
    Multilingual & 300 & May 2025 & Same pipeline, 9 non-Python
      languages~\cite{swebench2025multilingual} \\
    Pro & 1{,}865 & Sep 2025 & Long-horizon, multi-file; held-out proprietary
      repos~\cite{deng2025swebenchpro} \\
    \bottomrule
  \end{tabularx}
  \medskip
\end{frame}

\begin{frame}{SWE-bench: how hard is each task?}
  \label{fr:difficulty}
  \textbf{Verified tasks are devided into three categories}
  \medskip
  \begin{center}
    \small
    \begin{tabular}{@{}ll@{}}
      \toprule
      \textbf{Tier} & \textbf{Estimated human fix-time} \\
      \midrule
      \textcolor{TUMOrange}{$\blacksquare$}~Easy & $<$15\,min \\
      \textcolor{TUMGreen}{$\blacksquare$}~Medium & 15\,min\,--\,1\,h \\
      \textcolor{TUMExtRed}{$\blacksquare$}~Hard & $>$1\,h \\
      \bottomrule
    \end{tabular}
  \end{center}
\end{frame}

\begin{frame}{SWE-bench: the headline climbs}
  \label{fr:headline}
  \smallskip
  \textbf{SWE-bench Verified} --- nine systems, Apr'24 to
  Mar'25~\cite{ganhotra2025difficulty}.

  \smallskip
  \begin{center}
    \begin{tikzpicture}
      \begin{axis}[
        width=0.88\textwidth,
        height=5cm,
        ylabel={\% resolved},
        ymin=0, ymax=104,
        ytick={0,20,40,60,80,100},
        xmin=2024.16, xmax=2025.30,
        xtick={2024.25, 2024.5, 2024.75, 2025.0, 2025.25},
        xticklabels={Apr'24, Jul'24, Oct'24, Jan'25, Apr'25},
        xticklabel style={font=\scriptsize},
        yticklabel style={font=\scriptsize},
        ylabel style={font=\scriptsize},
        grid=major,
        grid style={dashed, TUMGrayLight},
        axis line style={TUMGrayDark},
        tick style={TUMGrayDark},
        legend style={font=\tiny, at={(0.5,1.02)}, anchor=south,
                      legend columns=4, draw=none, fill=none,
                      /tikz/every even column/.append style={column sep=4pt}},
        every axis plot/.append style={thick, mark size=1.6pt},
      ]
      \addplot[TUMBlue, mark=*] coordinates {
        (2024.251,18.2) (2024.467,33.6) (2024.571,23.2) (2024.918,50.8)
        (2024.918,55.0) (2024.975,57.2) (2025.044,64.6) (2025.099,63.4)
        (2025.203,65.4)
      };
      \addlegendentry{Overall}
      
      \addplot[TUMOrange, mark=square*] coordinates {
        (2024.251,27.3) (2024.467,47.9) (2024.571,36.6) (2024.918,70.6)
        (2024.918,72.2) (2024.975,74.7) (2025.044,77.3) (2025.099,81.4)
        (2025.203,80.4)
      };
      \addlegendentry{Easy}

      \addplot[TUMGreen, mark=triangle*] coordinates {
        (2024.251,10.0) (2024.467,28.0) (2024.571,16.9) (2024.918,42.5)
        (2024.918,49.4) (2024.975,50.6) (2025.044,62.1) (2025.099,56.7)
        (2025.203,62.1)
      };
      \addlegendentry{Medium}

      \addplot[TUMExtRed, mark=diamond*] coordinates {
        (2024.251,0.0) (2024.467,4.4) (2024.571,2.2) (2024.918,13.3)
        (2024.918,13.3) (2024.975,20.0) (2025.044,24.4) (2025.099,24.4)
        (2025.203,20.0)
      };
      \addlegendentry{Hard}
      \end{axis}
    \end{tikzpicture}
  \end{center}

  \smallskip
  \textbf{Easy near 80\,\%, hard stuck at 20\,\%.} The overall line hides the
  spread --- exact splits next.
\end{frame}

\begin{frame}{SWE-bench: what the numbers show}
  \label{fr:numbers}
  The same nine submissions, by annotator
  difficulty~\cite{ganhotra2025difficulty}:

  \medskip
  \begin{center}
    \begin{minipage}[c]{0.56\textwidth}
      \raggedleft\small
      \begin{tabular}{@{}lcc@{}}
        \toprule
        \textbf{Tier} & \textbf{First} & \textbf{Last} \\
        & {\scriptsize Apr '24} & {\scriptsize Mar '25} \\
        \midrule
        \textcolor{TUMOrange}{$\blacksquare$}~Easy \ {\small$<$15\,min
          ($n{=}194$)} & 27.3\,\% & 80.4\,\% \\
        \textcolor{TUMGreen}{$\blacksquare$}~Medium \ {\small$\leq$1\,h
          ($n{=}261$)} & 10.0\,\% & 62.1\,\% \\
        \textcolor{TUMExtRed}{$\blacksquare$}~Hard \ {\small$>$1\,h
          ($n{=}45$)} & \phantom{0}0.0\,\% & 20.0\,\% \\
        \bottomrule
      \end{tabular}
    \end{minipage}%
    \hspace{0.5cm}%
    \begin{minipage}[c]{0.30\textwidth}
      \begin{tikzpicture}
        \def\r{1.15cm}
        \fill[TUMOrange] (0,0) -- (90:\r)
          arc[start angle=90, end angle=-49.68, radius=\r] -- cycle;
        \fill[TUMGreen] (0,0) -- (-49.68:\r)
          arc[start angle=-49.68, end angle=-237.6, radius=\r] -- cycle;
        \fill[TUMExtRed] (0,0) -- (-237.6:\r)
          arc[start angle=-237.6, end angle=-270, radius=\r] -- cycle;
        \node[font=\scriptsize, text=white] at (20.16:0.62cm) {39\,\%};
        \node[font=\scriptsize, text=white] at (-143.64:0.62cm) {52\,\%};
        \node[font=\scriptsize, text=TUMExtRed] at (106.2:1.5cm) {9\,\%};
        \node[font=\scriptsize, text=TUMGrayDark, anchor=north]
          at (0,-1.35cm) {Task mix ($n{=}500$)};
      \end{tikzpicture}
    \end{minipage}
  \end{center}

  \smallskip
  \begin{itemize}
    \item Easy $+$ medium are \textbf{91\,\% of Verified}.
    \item The hard tier \textbf{stopped} near 20\,\%.
  \end{itemize}
\end{frame}

\section{What are the challenges?}

\begin{frame}{Challenge: intent capture}
  \label{fr:ambiguity}
  \begin{itemize}
    \item No system in the survey \textbf{asks a clarifying question}. It
      guesses and commits.
      \medskip
    \item Tests do not catch a wrong guess about something they never check.
  \end{itemize}

  \vfill
  {\footnotesize
   Clarification mechanisms exist and
   work~\cite{lahiri2022ticoder,mu2024clarifygpt,vijayvargiya2026ambiguity}.}
\end{frame}

\begin{frame}{Challenge: plan}
  \begin{itemize}
    \item Big repositories \textbf{do not fit in context}. The agent must
      guess what to read.
      \medskip
    \item Success drops from \textbf{80\,\%} on easy tasks to
      \textbf{20\,\%} on hard, multi-file ones (frame~\ref{fr:numbers}).
  \end{itemize}

  \vfill
  {\footnotesize
   Difficulty split from~\cite{ganhotra2025difficulty}; training-time fixes
   proposed by~\cite{wei2025swerl,pan2024swegym}.}
\end{frame}

\begin{frame}{Challenge: cost and verification}
  \begin{itemize}
    \item Running more attempts costs more but \textbf{does not raise the
      success rate}.
      \medskip
    \item \textbf{``All tests pass''} is not proof of correctness --- a wrong
      patch can still pass.
  \end{itemize}

  \vfill
  {\footnotesize
   Sample-and-rerank cost vs.\ resolve rate
   from~\cite{turncontrol2025,sweeffi2025}.}
\end{frame}

\begin{frame}{Takeaways}
  \begin{itemize}
    \item \textbf{Agentic beats one-shot because it can run things}, not
      because it reasons more.
      \medskip
    \item All three challenges are the \textbf{same gap}: nothing outside the
      model catches a wrong guess about intent, a wrong file, or a wrong fix.
      \medskip
    \item \textbf{Verification is the bottleneck.} In a randomised trial,
      developers were 19\,\% \emph{slower} --- believing they were 20\,\%
      faster~\cite{metr2025rct}.
  \end{itemize}

  \bigskip
  \begin{block}{}
    \centering
    Autocomplete to autonomous issue resolution: crossed.\\
    Impressive to \textbf{dependable}: still an open question.
  \end{block}
\end{frame}

\setbeamertemplate{footline}{}

\appendix

\begin{frame}[noframenumbering]{Backup: SWE-bench --- can we trust the numbers?}
  \label{fr:trust}
  \textbf{Verified against Pro}, April 2026:

  \medskip
  \begin{center}
    \small
    \begin{tabular}{@{}lccl@{}}
      \toprule
      & \textbf{Verified} & \textbf{Pro} & \textbf{Pro scored by} \\
      \midrule
      Claude Opus 4.5 & 80.9\,\% & 45.9\,\%
        & \textbf{standardised harness} \\
      GPT-5.2 & 80.0\,\% & 56.4\,\% & {\color{TUMGrayDark}\emph{self-reported}} \\
      GPT-5.3 & 85\phantom{.0}\,\% & 56.8\,\% & {\color{TUMGrayDark}\emph{self-reported}} \\
      GLM-5.1 & 78\phantom{.0}\,\% & 58.4\,\% & {\color{TUMGrayDark}\emph{self-reported}} \\
      \bottomrule
    \end{tabular}
  \end{center}

  \medskip
  {\footnotesize Read the \textbf{floor} only --- every model drops $\sim$20
  percentage points or more. \textbf{Not the spread:} the one measured row drops
  \emph{furthest}, so the ranking is an artefact of who
  scored~\cite{deng2025swebenchpro,agentmarketcap2026gap}.}

  \medskip
  \begin{itemize}
    \item \textbf{Contamination.} OpenAI dropped Verified in 2026 --- models
      reproduce gold patches from the task ID alone~\cite{openai2026swebench}.
    \item \textbf{Reward hacking.} Closing test leaks cuts scores by double
      digits: agents partly learn the grader~\cite{cursor2025rewardhacking}.
  \end{itemize}
\end{frame}

\begin{frame}[noframenumbering]{Backup: surveyed systems on SWE-bench Verified}
  {\scriptsize Source: \url{https://swebench.com}, ``All agents'' view.
  Best published result per system.}

  \medskip
  \small
  \begin{tabular}{@{}L{0.30\textwidth}L{0.30\textwidth}rl@{}}
    \toprule
    \textbf{System} & \textbf{Backbone} & \textbf{\% Resolved} & \textbf{Date} \\
    \midrule
    SAGE (OpenHands)
      & Claude 4.5 Opus   & 73.8\,\% & 2025-11 \\
    OpenHands
      & GPT-5             & 71.8\,\% & 2025-08 \\
    SWE-agent
      & Claude 4 Sonnet   & 66.6\,\% & 2025-05 \\
    AutoCodeRover-v2.1
      & Claude 3.5 Sonnet & 51.6\,\% & 2025-01 \\
    Agentless-1.5
      & Claude 3.5 Sonnet & 50.8\,\% & 2024-12 \\
    Agentless-1.5
      & GPT-4o            & 38.8\,\% & 2024-10 \\
    AutoCodeRover
      & GPT-4o            & 38.4\,\% & 2024-06 \\
    SWE-agent
      & GPT-4o            & 23.2\,\% & 2024-07 \\
    Devin
      & undisclosed       & 13.9\,\%$^*$ & 2024-03 \\
    \bottomrule
  \end{tabular}

  \smallskip
  {\scriptsize $^*$On \textbf{SWE-bench Full}, a different and easier split
  than the Verified numbers above --- not directly comparable. Self-reported;
  Cognition stopped publishing benchmark numbers in
  2024~\cite{cognition2024devin}.}
\end{frame}

\begin{frame}[noframenumbering]{Backup: the context bottleneck}
  \begin{itemize}
    \item Longer windows do not settle \emph{what to read}. Keeping full
      dependency context works best; padding with loosely related code
      \textbf{hurts} correctness~\cite{repoexec2025}.
      \medskip
    \item Appending every observation causes context explosion and semantic
      drift: early critical information gets buried~\cite{contextastool2025}.
      \medskip
    \item The whole prefix is re-sent each turn, so tokens over $n$ turns grow as
      $O(n^2)$~\cite{turncontrol2025}. A protocol property, not a finding;
      caching cuts the constant, not the growth.
  \end{itemize}

  \bigskip
  {\small An agent cannot fix what it never places in context.}
\end{frame}

\begin{frame}[noframenumbering]{Backup: Plan --- localisation}
  \label{fr:localisation}
  {\large So how does it \textbf{stop guessing}?}

  \medskip
  Point at real code. Narrow the repo from broad to specific --- each step
  something you can check:

  \bigskip
  \begin{itemize}
    \item \textbf{Agentless} --- files $\rightarrow$ classes/functions
      $\rightarrow$ edit locations~\cite{xia2024agentless}
      \medskip
    \item \textbf{AutoCodeRover} --- AST-backed search APIs, coarse to
      fine~\cite{zhang2024autocoderover}
  \end{itemize}

  \bigskip
  \begin{block}{What the useful ones share}
    They name \textbf{concrete program elements} that can be checked against the
    repository --- not invented requirements.
  \end{block}
\end{frame}

\begin{frame}[noframenumbering]{Backup: Localisation --- Agentless}
  {\small Narrowing without letting the model wander: three fixed steps, coarse
  to fine~\cite{xia2024agentless}.}

  \bigskip
  \begin{enumerate}
    \item \textbf{Rank files} --- from the repo layout and the issue, shortlist
      the files most likely involved.
      \medskip
    \item \textbf{Rank elements} --- inside those files, narrow to specific
      classes and functions.
      \medskip
    \item \textbf{Pin locations} --- inside those, the exact lines to edit.
  \end{enumerate}

  \bigskip
  \begin{block}{Fixed, not discovered}
    Three prompts, coarse to fine. No search, no tools --- the narrowing is
    hard-coded into the scaffold.
  \end{block}
\end{frame}

\begin{frame}[noframenumbering]{Backup: Localisation --- AutoCodeRover}
  {\small Let the model navigate --- but through the syntax tree, not by reading
  whole files~\cite{zhang2024autocoderover}.}

  \bigskip
  \begin{itemize}
    \item \textbf{Search APIs} --- \texttt{search\_class},
      \texttt{search\_method}, \texttt{search\_method\_in\_class}.
      \medskip
    \item The agent calls them \textbf{coarse to fine}, pulling in just the code
      it needs.
      \medskip
    \item AST-backed: structured, cheaper context than dumping whole files into
      the prompt.
  \end{itemize}

  \bigskip
  \begin{block}{Discovered, not fixed}
    The model chooses what to search --- agentic, where Agentless is a fixed
    funnel.
  \end{block}
\end{frame}

\setbeamertemplate{bibliography item}[text]
\setbeamerfont{bibliography item}{size=\tiny}
\setbeamerfont{bibliography entry author}{size=\tiny}
\setbeamerfont{bibliography entry title}{size=\tiny}
\setbeamerfont{bibliography entry location}{size=\tiny}
\setbeamerfont{bibliography entry note}{size=\tiny}

\setbeamertemplate{bibliography entry title}{}
\setbeamertemplate{bibliography entry location}{}
\setbeamertemplate{bibliography entry note}{}

\makeatletter
\let\beamer@oldthebibliography\thebibliography
\renewcommand{\thebibliography}[1]{%
  \beamer@oldthebibliography{#1}%
  \setlength{\itemsep}{0pt}%
  \setlength{\parsep}{0pt}%
  \setlength{\topsep}{0pt}%
}
\makeatother

\begin{frame}[noframenumbering, allowframebreaks]{References}
  \tiny
  \bibliographystyle{abbrv}
  \bibliography{references,slides-extra}
\end{frame}

\end{document}
